\documentclass{article}
\usepackage[english]{babel}

\usepackage[a4paper,top=2cm,bottom=2cm,left=3cm,right=3cm,marginparwidth=1.75cm]{geometry}
\usepackage{amsmath}
\usepackage{graphicx}
\usepackage[colorlinks=true, allcolors=blue]{hyperref}
\usepackage{listings}
\usepackage{minted}
\usepackage{booktabs}
\setlength{\parindent}{0pt}

\title{Physics of Complex Systems: Exercise }
\author{Leopold Schaller and Daniel Lender}

\begin{document}
\maketitle

\section*{Task 2.1: Drunk Drivers}
\subsection*{(a)}
The second table displays the conditional probability
\begin{align}
    P \left( female \mid age \right)
\end{align}
an we can calculate the total probability with summing over all age intervalls. We obtain
\begin{align}
    P \left( female \right) &= \sum_{ age} P \left( total \right) \cdot P \left( female \mid age \right)\\
                              &= 12.14 \%
\end{align}

\subsection*{(b)}
    Bayes formula is
\begin{align}
    P \left( A \mid B \right) = P \left( B \mid A \right) \cdot \frac{P \left( A  \right)}{P \left( B \right)}
\end{align}
    and is used to exchanges the conditions in conditional probabilities. For our example, we have to use the formula:
\begin{align}
    P \left( age \mid female \right) = P \left( female \mid age \right) \cdot \frac{P \left( age  \right)}{P \left( female \right)}
\end{align}
    We can calculate this now by using the values from the first table and the calculated value $P \left( female \right)$ and obtain the folowing table:
\\
\begin{table}[h]
\centering
\begin{tabular}{l|ccccccc}
\toprule
Age group & 18--24 & 25--34 & 35--44 & 35--54 & 55--64 & $>65$ & Sum \\
\midrule
$P(A)$ & 26.1 & 22.4 & 18.9 & 18.5 & 8.7 & 5.4 & 100 \\
$P(B \mid A)$ & 9.1 & 11.9 & 15.2 & 14.4 & 12.3 & 9.2 & -- \\
$P(A \mid B)$ & 19.6 & 21.9 & 23.7 & 21.9 & 8.8 & 4.1 & 100 \\
\bottomrule
\end{tabular}
\end{table}

\section*{Task 2.2: Poisson distribution}
We start with the bonomial distribution
\[
P(X = k) = \binom{N}{k} p^k (1 - p)^{N-k}
\]
and want to calculate the Poisson distribution.
\subsection*{(a)}
The term $(1 - p)^{N-k}$ is approximated as follows:
\begin{align}
\lim_{N \to \infty} (1 - p)^{N-k} &= \lim_{N \to \infty}
\left[
\cdot \left(1 - \frac{\lambda}{N}\right)^N
\cdot \left(1 - \frac{\lambda}{N}\right)^{-k}
\right] \\
\\
&=
\cdot \underbrace{\lim_{N \to \infty} \left(1 - \frac{\lambda}{N}\right)^N}_{=e^{-\lambda}}
\cdot \underbrace{\lim_{N \to \infty} \left(1 - \frac{\lambda}{N}\right)^{-k}}_{=1} \\
\\
&= e^{-\lambda}
\end{align}

\subsection*{(b)}
The term $\binom{N}{k} p^k$ is approximated using the definition of the binomial coefficient:
\begin{align}
\lim_{N \to \infty} \frac{N!}{k!(N-k)!}
\cdot \frac{\lambda^k}{N^k}
&= \lim_{N \to \infty}
\left[
\frac{N!}{(N-k)! \, N^k}
\cdot \frac{\lambda^k}{k!}
\right] \\
\\
&= \frac{\lambda^k}{k!}
\cdot \underbrace{\lim_{N \to \infty} \frac{N!}{(N-k)! \, N^k}}_{=1}
\\
&= \frac{\lambda^k}{k!}
\end{align}

\subsection*{(c)}
    If we combine the results of (a) and (b), we obtain the Poisson distribution:
    \begin{align}
      \lim_{N \to \infty} \underbrace{\binom{N}{k} p^k_}_{\frac{\lambda^k}{k!}} \underbrace{(1 - p)^{N-k}}_{e^{-\lambda}} = \frac{\lambda^k}{k!} \cdot  e^{-\lambda}
    \end{align}

\subsection*{(d)}
For normalization, we have to perform the sum over all events $k$ and should obtain a value of $1$. Using the series expansion of the exponential function, we get the normalization condition. This is shown in the following:
\begin{align}
\sum_{k=0}^{\infty} P(X = k)
&= \sum_{k=0}^{\infty} \frac{\lambda^k}{k!} e^{-\lambda} \\
&= e^{-\lambda} \sum_{k=0}^{\infty} \frac{\lambda^k}{k!} \\
&= e^{-\lambda} \cdot e^{\lambda} \\
&= 1
\end{align}

\subsection*{(e)}
The moment generating $M_X(t)$ function is the expectation value of $e^{tX}$, where $X$ is the random variable. In our case, this variable follows the Poisson distribution.
\begin{align}
M_X(t)
&= E(e^{tX}) \\
&= \sum_{k=0}^{\infty} e^{tk} P(X=k) \\
&= \sum_{k=0}^{\infty} e^{tk} \frac{\lambda^k}{k!} e^{-\lambda} \\
&= e^{-\lambda} \sum_{k=0}^{\infty} \frac{(\lambda e^t)^k}{k!} \\
&= e^{-\lambda} e^{\lambda e^t} \\
&= e^{\lambda(e^t-1)}.
\end{align}

The cumulant-generating function $C_X(t)$ is the logarithm of $M_X(t)$, so it is simply
\begin{align}
    C_X(t) =\lambda \left( e^t - 1 \right) \text{ .}
\end{align}

\subsection*{(f)}
If we want to determine an cumulant $\kappa_n$ of n'th order, we have to use the following relation.
\begin{align}
\kappa_n
&= \frac{d^n}{dt^n} C_X(t) \right|_{t=0} \\
&= \frac{d^n}{dt^n} \lambda \left( e^t - 1 \right) \right|_{t=0} \\
&= \frac{d^n}{dt^n} \lambda e^t  \right|_{t=0} \\
&= \lambda
\end{align}
We obtain, that for the Poisson distribution all cumulants are the same. This means that especially the expectation value $E(X)$ and the variance $Var(X)$ are the same.


\section*{Task 2.3: Reconstruct probabilities from given moments}
\subsection*{(a)}
We have three unknown variables, thus we need three equations to determine the variables (here probabilities) unequivocally. The first equation is the normalization boundary, for the second and the third equation one moment each is needed. Therefore we need two moments to determine $P_a$, $P_b$ and $P_c$.

\subsection*{(b)}
\begin{align}
\left.
\begin{aligned}
&P_a+\ \ P_b+\ \ \ \ P_c=\ 1\\
&P_a+3P_b+5\ \ P_c=M_1\\
&P_a+9P_b+25P_c=M_2
\end{aligned}
\right\}
\Leftrightarrow \begin{pmatrix}
1&1&1\\
1&3&5\\
1&9&25
\end{pmatrix}
\cdot
\begin{pmatrix}
P_a\\
P_b\\
P_c
\end{pmatrix}
=
\begin{pmatrix}
1\\
M_1\\
M_2
\end{pmatrix}
\end{align}
We have used the following code to solve this system of equations:
\begin{minted}{python}
import sympy as sp
A = sp.Matrix([
    [1, 1, 1],
    [1, 3, 5],
    [1, 9, 25]
])

b1, M1, M2 = sp.symbols('b1 M1 M2')
B = sp.Matrix([b1, M1, M2])

# 3. Solve the system Ax = B
x = A.solve(B)

numeric_x = x.subs({b1: 1})
sp.pprint(numeric_x)
\end{minted}
This results in the following formulas for the probabilities:
\begin{align}
    P_a &= -M_1 + \frac{M_2}{8} + \frac{15}{8}\\
    P_b &= 3 \frac{M_1}{2} - \frac{M_2}{4} + \frac{5}{4}\\
    P_c &= - \frac{M_1}{2} - \frac{M_2}{8} + \frac{3}{8}
\end{align}

\subsection*{(c)}
The moments $M_n$ are strictly connected to the map $X$ and the probabilities $P_i$ via
\begin{align}
    M_n = \sum_{i} X_{i}^{n} \cdot P_i
\end{align}
The probabilities $P_i$ are not specified as positive numbers, but only that their sum has to be equal to one. Therefore, we can construct moments $M_1$ and $M_2$ which will lead to negative probabilities. This would not be a consistent result. For example, we can use
\begin{align}
    M_1 &= 10 \\
    M_2 &= 25
\end{align}
and will get probabilities $P_a = -5$, $P_b = 7.5$ and $P_c = -2.5$. They sum up to 1, but are not probabilities.

\subsection*{(d)}
The moments
\begin{align}
M_1 &= 3.5 \\
M_2 &= 15
\end{align}
lead to the consistent probabilities
\begin{align}
P_a &= 0.25 \\
P_b &= 0.25 \\
P_c &= 0.5
\end{align}
The formula for the nth moment is now
\begin{align}
M_n = 0.25 \cdot 1^n + 0.25 \cdot 3^n + 0.5 \cdot 5^n
\end{align}
and the first moments are:
\begin{table}[h]
\centering
\begin{tabular}{l|ccccccc}
\toprule
n & 0 & 1 & 2 & 3 & 4 & 5 \\
\midrule
$M_n$ & 1 & 3.5 & 15 & 69.5 & 333 & 1623 \\
\bottomrule
\end{tabular}
\end{table}

\end{document}

# latexmk -pdf -shell-escape -output-directory=output Uebung2/exercise2/exercise2_Schaller_Lender.tex
